% HD-N25. GPT-6 Astra Pro /root/imprimitive, 21 November 2026.
% Inherited effort; numerical setting not exposed.
\section{HD-N25: every two-level positive graph coupling}
\label{sec:two-level-positive-graph-cones}

Take any nodal arrangement of the actual $N=9408$ translated
abelian threefolds $C_\nu\subset A$ in HD-N23. Partition its
components into two sets, and put
\[
 P=\bigoplus_{\nu\in I}\mathcal O_{C_\nu},\qquad
 Q=\bigoplus_{\nu\in J}\mathcal O_{C_\nu},\qquad
 I\amalg J=\{1,\ldots,N\}.
                                                               \tag{N25.1}
\]
For every actual morphism
$\delta:P\to Q[3]$, let
\[
             Q[2]\xrightarrow{i}E_\delta\xrightarrow{q}P
                     \xrightarrow{\delta}Q[3]              \tag{N25.2}
\]
be its fiber triangle. This is an actual perfect gluable
complex, with whole character
\[
             \operatorname{ch}(E_\delta)=32\theta^3+120T.
                                                               \tag{N25.3}
\]
Choose a partition placing the graph directions $(1,1,1)$
and $(i,i,-1)$ in different sets. Their sixteen intersection
classes give nonzero choices of $\delta$, realized by the top
Koszul class at an actual transverse intersection point.
Thus the construction includes nontrivial coupled objects
with positive cycle signs.

Nevertheless, for every $\delta$, every actual coherent action
of a faithful finite point subgroup on $E_\delta$ has a
nonzero weak degree-two trace kernel of dimension at least
$495$. The same statement holds for every actual isogeny
pullback of $E_\delta$, including newly available flat-tensor
point stabilizers, followed by actual line tensor and
Fourier--Mukai equivalences. The proof uses the canonical two-cohomology
filtration and does not assume the complex is simple, the
coupling descends to a quotient, or a chosen direct-sum
projector lifts to an endomorphism of $E_\delta$.

\paragraph{1. The exact cross Ext and self-Ext maps.}
Two distinct components of the nodal arrangement are either
disjoint or meet transversely in finitely many points. At a
transverse point the actual local ring and modules are
\[
 R=\mathbf C[[x_1,x_2,x_3,y_1,y_2,y_3]],\qquad
 U=R/(x),\qquad V=R/(y).
\]
Applying the $x$-Koszul resolution of $U$ to $V$ gives the
cochain Koszul complex of a regular sequence of length three.
It has cohomology only in degree three, where it is the
one-dimensional residue-field module. The reverse calculation
is identical. Since a sheaf with finite support has no positive
cohomology, the actual local-to-global maps prove
\[
 \operatorname{Ext}^k(\mathcal O_{C_\mu},\mathcal O_{C_\nu})=0
       \quad(\mu\ne\nu,\ k\ne3).                         \tag{N25.4}
\]
The degree-three group is the direct sum of its one-dimensional
local Koszul groups over their intersection points. This
supplies the asserted actual nonzero couplings; for example,
the two graph directions $(1,1,1)$ and $(i,i,-1)$ meet in
sixteen points for any of their translates.

For one component $C=B+a$, its self-Ext groups have dimensions
\[
             \dim\operatorname{Ext}^k(\mathcal O_C,\mathcal O_C)
                                   =\binom6k.              \tag{N25.5}
\]
Here is an actual product map giving that calculation and the
point-action statement needed below. Choose an ample line
bundle $L$ on the ambient abelian sixfold $X$, let
$j:B\hookrightarrow X$, and put
\[
 D=\ker(j^\vee\lambda_L:X\to\widehat B)^0,
 \qquad\pi:B\times D\to X,\quad(b,d)\mapsto b+d.           \tag{N25.6}
\]
The restriction of $j^\vee\lambda_L$ to $B$ is the
polarization isogeny of $L|_B$, so $B\cap D$ is finite;
this proves that the displayed addition map is an isogeny.
The exact primary locator, including this kernel construction,
is Milne, \emph{Abelian Varieties}, version2.00, ChapterI,
Proposition10.1 and its proof, printedp.42, as checked in HD-N21.
The inverse image of $C$ is a disjoint union of $B\times\{d\}$
after a translation, and the restriction of $\pi$ to each
component is a translation isomorphism onto $C$.
Finite flat base change gives its literal pullback sheaves.

On a component the derived Kunneth map is
\[
 R\operatorname{Hom}_{B\times D}
 (\mathcal O_B\boxtimes\mathcal O_d,
  \mathcal O_B\boxtimes\mathcal O_d)
 =R\Gamma(B,\mathcal O_B)\otimes
                         R\operatorname{Hom}_D(\mathcal O_d,\mathcal O_d).
                                                               \tag{N25.7}
\]
Invariant Dolbeault forms on $B$ and the actual point Koszul
resolution on $D$ give exterior cohomologies on respectively
$H^1(B,\mathcal O_B)$ and $T_dD$, each of dimension three.
The deck group of $\pi$ permutes its inverse-image components
freely and transitively. Taking its canonical descent
invariants identifies the self-Ext downstairs with one copy
of this exterior algebra, proving (N25.5). One may check
that descent map directly by the adjunction unit and averaging:
$(\pi_*\mathcal O_{B\times D})^{\ker\pi}=\mathcal O_X$;
invariants are exact in characteristic zero.

The same product map proves the following precise stabilizer
fact. Suppose $T_b^*(-)\otimes M$ carries $\mathcal O_C$ to
itself, where $M\in\operatorname{Pic}^0(X)$. Then $b\in B$
and $M|_B$ is trivial. Lift the translation to $(b,0)$ in
(N25.6). The actual flat-line pullback has the form
$\pi^*M=\mathcal O_B\boxtimes M_D$. On (N25.7), translation
acts trivially on the invariant Dolbeault forms, while tensor
by $M_D$ acts trivially on the point self-Ext: a local line
trivialization multiplies its two arguments by the same
invertible scalar and cancels on the point Koszul Hom complex.
Changing that trivialization gives the same map on Ext.
Pullback on Ext is injective, split by the finite-etale trace
divided by $\deg\pi$. Therefore this point stabilizer acts
as the identity on every self-Ext group downstairs.
This proves the assertion also for nontrivial ambient flat
lines whose restriction to $B$ is trivial.

\paragraph{2. The canonical two-cohomology filtration.}
Equation(N25.4) gives
\[
             \operatorname{Ext}^k(P,Q)=
             \operatorname{Ext}^k(Q,P)=0\quad(k\ne3).
                                                               \tag{N25.8}
\]
Taking cohomology in (N25.2) identifies
\[
          \mathcal H^{-2}(E_\delta)=Q,\qquad
          \mathcal H^0(E_\delta)=P,
          \qquad\mathcal H^j(E_\delta)=0\ (j\ne-2,0).
                                                               \tag{N25.9}
\]
Thus (N25.2) is its canonical Postnikov triangle, with the
specified connecting morphism. In particular every t-exact
point functor and every isomorphism of $E_\delta$ preserve
this triangle functorially. No ordinary decomposition of the
whole complex is being assumed.

The complex is gluable. For $k<0$, apply Hom from $E_\delta$
to the shifted triangle. The two possible adjacent groups
$\operatorname{Hom}(E_\delta,P[k])$ and
$\operatorname{Hom}(E_\delta,Q[k+2])$ vanish: applying Hom
into these targets from the unshifted triangle reduces them
to negative self-Ext of $P,Q$ and the cross groups
$\operatorname{Ext}^{k-2}(Q,P)$ and
$\operatorname{Ext}^{k+2}(P,Q)$, which vanish by (N25.8).
Hence $\operatorname{Ext}^{<0}(E_\delta,E_\delta)=0$.
Perfection follows from that of the smooth closed-support
sheaves and closure under the actual cone. Its character
(N25.3) follows because the two shifts are even and the
abelian embeddings have Todd class one, as in HD-N23.

\paragraph{3. The actual degree-two short exact sequence.}
The triangle gives canonical identifications
\[
 \begin{aligned}
 \operatorname{Hom}(E_\delta,Q[4])
        &\xrightarrow[\ i^*\ ]{\sim}\operatorname{Ext}^2(Q,Q),\\
 \operatorname{Ext}^j(P,P)
        &\xrightarrow[\ q^*\ ]{\sim}\operatorname{Hom}(E_\delta,P[j])
                         \quad(j=1,2).
 \end{aligned}                                               \tag{N25.10}
\]
For the first, the neighboring cross groups are
$\operatorname{Ext}^4(P,Q)$ and $\operatorname{Ext}^5(P,Q)$.
For the second, they are the appropriate negative or degree-zero
cross groups from $Q$ to $P$. All vanish by (N25.8).
For $a\in\operatorname{Ext}^2(Q,Q)$, take the unique
$u:E_\delta\to Q[4]$ restricting to $a[2]$, and send it to
$i[2]u$. For $\xi\in\operatorname{Ext}^2(E_\delta,E_\delta)$,
let $b\in\operatorname{Ext}^2(P,P)$ be the unique map with
$q[2]\xi=bq$. These are the actual maps in
\[
 0\longrightarrow\operatorname{Ext}^2(Q,Q)
 \longrightarrow\operatorname{Ext}^2(E_\delta,E_\delta)
 \longrightarrow\operatorname{Ext}^2(P,P)\longrightarrow0.
                                                               \tag{N25.11}
\]
We check both possible boundary maps explicitly. Any
$E_\delta\to P[1]$ is $cq$ for $c:P\to P[1]$;
its boundary in $\operatorname{Hom}(E_\delta,Q[4])$ is
$\delta[1]cq$, and the middle composite belongs to
$\operatorname{Ext}^4(P,Q)=0$. Thus the first map in
(N25.11) is injective. The obstruction to lifting $b$ from
the right is $\delta[2]bq$; its middle composite belongs to
$\operatorname{Ext}^5(P,Q)=0$, so it vanishes. This proves
surjectivity. Exactness in the middle is the same Hom long
exact sequence. Every construction used (N25.2), so the
sequence is natural for all coherent point actions.

The distinct-support cross groups contribute nothing in
degree two. Equations(N25.4)--(N25.5) therefore give the exact
ordinary dimension
\[
                  \dim\operatorname{Ext}^2(E_\delta,E_\delta)
                                  =15N=141120              \tag{N25.12}
\]
for every $\delta$. No degeneration of an endomorphism
spectral sequence is used.

\paragraph{4. Every coherent point action on the original arrangement.}
Let a point functor $T_a^*(-)\otimes M$ preserve $E_\delta$.
It preserves its two cohomology sheaves (N25.9) separately.
Maps between structure sheaves on distinct supports vanish
in degree zero by (N25.4); each individual support sheaf has
scalar endomorphisms. Hence the induced cohomology isomorphisms
permute the actual component supports and have only scalar
freedom on each supported line. They cannot mix distinct
supports by an arbitrary matrix.

It follows that $a$ preserves the whole union $Z$, and that
$M$ restricts trivially to every component direction. The
coordinate-threefold directions occurring in HD-N23 contain
all six coordinate elliptic factors between them. Under the
actual product identification
$\operatorname{Pic}^0(E_i^6)=\widehat E_i^6$, those restrictions
therefore force every coordinate of $M$ to be zero. Thus
$M=\mathcal O_A$. The translation part belongs to the full
$H=\operatorname{Stab}(Z)$ of HD-N23, with $|H|\le144$.
The coupling may shrink this stabilizer, which only improves
the following bound.

For any actual coherent faithful point group $G$ on the
complex, (N25.11) is an exact sequence of genuine $G$-modules:
any scalar coherence cancels on conjugation of self-Ext.
Taking invariants is exact. On the two exterior self-Ext sums,
the group permutes component blocks; each block stabilizer
acts as the identity by the actual complement argument in
paragraph1. Consequently each component orbit contributes
exactly fifteen invariant dimensions. This statement is
independent of all scalar choices of the cohomology
isomorphisms. The two grades are preserved, and their total
number of component orbits is at least
$\lceil9408/144\rceil=66$. Therefore
\[
         \dim\operatorname{Ext}^2(E_\delta,E_\delta)^G\ge990.
                                                               \tag{N25.13}
\]
The ambient full degree-two trace target has dimension495.
Thus its invariant kernel has dimension at least495. The
identity-sector comparison of HD-N12 gives the same weak
categorical trace kernel for the actual equivariant object.
No simplicity hypothesis or invariant endomorphism projector
of the whole complex is required.

\paragraph{5. All point actions after isogeny pullback.}
Let $f:X\to A$ be an actual isogeny. Its inverse images of
the $C_\nu$ are disjoint unions of translated connected abelian
threefolds. They still meet only in transverse double points;
no triple point is introduced by an etale map. Pulling back
(N25.2) gives the same actual two-cohomology construction on
these components, so (N25.8)--(N25.11) hold there verbatim.

Let $G$ now be any actual coherent finite faithful point group
preserving $f^*E_\delta$. Its translation coordinate $a$ sends
the inverse-image support union to itself. Applying $f$ shows
that $f(a)\in H$. Its flat coordinate need not be trivial on
$X$, but the actual cohomology isomorphisms force its restriction
to every inverse-image component to be trivial. The complement
isogeny argument in paragraph1 applies on $X$ itself and
proves that every component stabilizer, including its flat
tensor coordinate, acts trivially on the component self-Ext.

Map each inverse-image component to its original component
label. This is a surjective map of finite sets and is equivariant
for the homomorphism $G\to H$ given by $a\mapsto f(a)$.
Thus the number of component orbits upstairs is at least
the number of orbits of the image in $H$ downstairs, hence
at least66. It is unnecessary to divide the new component
count by the possibly much larger isogeny degree. The exact
invariant sequence (N25.11) again gives (N25.13), and hence
the trace-kernel bound495, for all these coherent point actions.
This also deals with nonsimple cones and new pure-tensor
stabilizers without claiming descent of $\delta$.

Finally actual line tensor and Fourier--Mukai equivalences
transport self-Ext, trace and point actions by the precise
intertwiners proved in HD-N12 and HD-N17. A group preserving
the transported object conjugates back to a faithful point
group of the object just handled. Thus the same bound holds
through the listed transports. Arbitrary longer complexes,
new same-grade multiplicities, or different support arrangements
are not asserted to satisfy this particular bound.

\paragraph{Check scope.}
The universal proof uses the actual cross-Koszul complexes,
the complement addition isogeny, and four Hom exact-sequence
maps, with their two explicit zero factorizations. The separate
certificate \path{two_level_graph_ext_certificate.py} records
the exact exterior dimensions, all neighboring cross-degree
vanishings used above, the component-orbit integer inequalities,
and the resulting trace bounds. Those finite checks supplement,
and do not replace, the displayed equivariant morphisms.
